\section{不确定性感知}

\begin{frame}{声磁协同：连续递推与绝对约束互补}
  \begin{columns}[T,onlytextwidth]
    \begin{column}{0.30\textwidth}
      \begin{block}{高频骨架}
        IMU 约 100 Hz\\
        DVL 约 2 Hz\\[1mm]
        ES-EKF 保持连续状态
      \end{block}
    \end{column}
    \begin{column}{0.32\textwidth}
      \begin{block}{声学辅助}
        裸露、水质好：声呐主导\\
        掩埋、浑浊：质量分数下降
      \end{block}
    \end{column}
    \begin{column}{0.32\textwidth}
      \begin{block}{磁学拉回}
        DLIA 输出幅值、相位、SNR\\
        FWHM 提供垂直距离约束
      \end{block}
    \end{column}
  \end{columns}

  \vspace{4mm}
  \begin{beamercolorbox}[rounded=true,sep=2mm]{block body}
    \centering
    \[
      z_{\mathrm{mag}}
      = \frac{1}{2}v\sin\alpha\,\Delta t_{\mathrm{FWHM}},
      \qquad
      d_{\mathrm{depth}} = z_{\mathrm{mag}}-h_{\mathrm{alt}}
    \]
  \end{beamercolorbox}

  \begin{itemize}
    \item 几何时间跨度替代未知电流幅值，降低铠装衰减和姿态变化影响。
    \item SNR 与拟合残差进入 \(\boldsymbol{R}_k\)，决定磁观测应被信任到什么程度。
  \end{itemize}
\end{frame}

\begin{frame}{不确定性成为跨层调度信号}
  \centering
  \includegraphics[width=0.98\linewidth,height=0.86\textheight,keepaspectratio]
    {architecture/auv_v2_uncertainty_highway.pdf}
\end{frame}

\begin{frame}{状态估计证据：稳定输出且协方差机制被真实触发}
  \begin{columns}[c,onlytextwidth]
    \begin{column}{0.58\textwidth}
      \centering
      \includegraphics[width=\linewidth,height=0.61\textheight,keepaspectratio]
        {experiments/mag_lever_arm_fullflow_20260705_2145/02_mag_extrinsics_error_reduction.png}
      \scriptsize 仿真外参标定框架，不等同于真机转台标定。
    \end{column}
    \begin{column}{0.38\textwidth}
      \begin{block}{扰动回归}
        \centering
        \KeyNumber{24/24}{8 场景 \(\times\) 3 随机种子全部完成}
      \end{block}
      \begin{block}{自适应协方差}
        \centering
        \KeyNumber{0.24--0.36}{八类场景中的触发率范围}
      \end{block}
      \begin{block}{磁外参仿真标定}
        \centering
        \KeyNumber{96.27\% / 95.31\%}{平移误差 / 安装角误差下降}
      \end{block}
    \end{column}
  \end{columns}

  \FrameConclusion{“输出没有崩溃”与“协方差知道自己在变差”是两条独立证据。}
\end{frame}
